\section{Introduction}
\label{sec:introduction}

Connected and automated vehicles (CAVs) and connected human-driven vehicles
need information that may be visible only to other vehicles or roadside
infrastructure. Examples include an occluded road user, the current traffic
signal state, and a cooperative maneuver plan. Vehicle-to-everything (V2X) and
fifth-generation (5G) communication can deliver this information to automated
driving systems and human drivers
~\cite{etsi122186r19,one6g2025overview}.

These applications create two communication patterns. Connected-vehicle (CV)
and intelligent transportation system (ITS) applications commonly exchange
compact, stateless SAE International J2735 updates. A later update replaces an
earlier one without preserving an application session. Complex CAV
applications instead use correlated, stateful operations for cooperative
perception, trajectory planning, maneuver coordination, or remote assistance.
An event or fault can trigger an uplink-heavy or downlink-heavy burst, while
remote teleoperation can sustain video and control traffic. Both application
classes are proposed for direct V2X and network-assisted 5G communication, yet
they lack one application protocol and a common evaluation of their
communication limits. We therefore ask: \emph{Can today's V2X and 5G paths
support both CV and CAV applications, and at what workload size and operating
conditions does each path stop meeting the required response latency,
completion rate, or application goodput?}

We answer this question with the Intersection Programming Interface (IPI) and
a controlled, real-vehicle field study. IPI provides one transport-independent
application protocol for stateless CV/ITS messages and stateful CAV
operations. The study applies both workload classes to a commercial Long-Term
Evolution cellular V2X (LTE C-V2X) direct PC5 path and a private 5G New Radio
(NR) Uu path. The private network contains one physical CAV user equipment
(UE), a dedicated radio, and a clean 40-MHz n48 channel. Experimental traffic
introduces all contention. Across payload, direction, transport, radio
condition, configuration, competing traffic, demand, and interruption, we
measure response round-trip time (RTT), the percentage of issued requests
answered within each application deadline, and directional application
goodput.

This paper makes two contributions:
\begin{itemize}
  \item \textbf{A common CV, ITS, and CAV application protocol.} IPI carries
  typed or pre-encoded J2735 messages and correlated CAV operations through
  the same transport-independent interface.
  \item \textbf{A controlled communication performance evaluation.} A real
  autonomous vehicle carries representative CV and CAV workloads over direct
  PC5 and network-assisted Uu. The experiments locate the conditions under
  which each path misses application deadlines or required directional rates.
\end{itemize}

The measurements yield three insights:
\begin{enumerate}
  \renewcommand{\labelenumi}{\arabic{enumi})}
  \item \textbf{The measured commercial direct V2X path supports compact J2735
  messages only within a limited payload and coverage envelope.} Roadside-unit
  placement or a supplemental path is necessary where buildings, highway
  geometry, or sparse infrastructure leave weak or absent coverage.
  \item \textbf{The CAV application-packet size that the measured 5G uplink can
  complete within a decision deadline is severely limited.} CAV developers
  must select, compress, or progressively transmit vehicle-originated state
  until the uplink envelope expands.
  \item \textbf{5G/6G systems must support event-triggered CAV bursts,
  sustained directional streams, and concurrent deadline-sensitive traffic.}
  Radio configurations and schedulers must preserve both directions and
  isolate simultaneous flows according to application deadline and duration.
\end{enumerate}
